\begin{literature}[\bf Tensor Networks]
  The modern theory of \emph{tensor networks} shaped out of the study of \emph{renormalization}
  and quantum many-body systems, initially with the \emph{density-matrix-renormalization-group} of \textcite{whiteDensityMatrixFormulation1992,whiteDensitymatrixAlgorithmsQuantum1993} and later extended by \textcite{ostlundThermodynamicLimitDensity1995} as a variational method (ansatz?) in the form of \emph{matrix product states} (MPS) (Maybe this is actually due to \cite{fannesFinitelyCorrelatedStates1992}). This was further generalized to 2 dimensions with the introduction of \emph{projected entangled-pair states} (PEPS) \cite{verstraeteRenormalizationAlgorithmsQuantumMany2004}. \emph{Multiscale-entanglement renormalization ansatz} (MERA) \cite{vidalClassQuantumManybody2008} can be considered as an extension of MPS allowing for long-range entanglement (this is probably wrong). Whats missing? Holography \cite{singhHolographicSpinNetworks2018} but there are earlier ones.

  In the purely mathematical view as a network of tensor contractions however, both concept and the notation we use today were introduced by \cite{Penrose1971Applications} under the name of \emph{tensor-systems}.
\end{literature}

While we introduce this concept classically and endow it with symmetry as a start, we will start concerning ourselves with a categorical embedding of tensor networks in what is commonly reffereed to as a \emph{fusion-category} (Lit ???).

\begin{literature}[\bf Symmetric Tensor Network]
  Originally, the idea to exploit symmetries of the hamiltonian for a speedup was proposed to be used in dmrg by \cite{mccullochNonAbelianDensityMatrix2002} (What is \cite{sierraDensityMatrixRenormalization1998a}? This sounds like it as well.). More general theory for \emph{globally invariant tensor networks}, i.e. networks that are built of tensors that are globally invariant is given by \textcite{singhTensorNetworkDecompositions2010a} preceeding more detailed accounts on abelian symmetries \cite{singhTensorNetworkStates2011} and nonabelian ($\SU{2}$) symmetries in \cite{singhTensorNetworkStates2012a}. A similar approach to (non)abelian symmetries is given by \textcite{weichselbaumNonabelianSymmetriesTensor2012, weichselbaumXsymbolsNonAbelianSymmetries2020}. These rely on the general idea, that under representations of some compact group $G$, there are on site symmetries on each tensor, i.e. $[A, U_g^{\otimes k}] = 0$ there is a generalized \emph{wigner-eckart-theorem} for tensors with $n$ legs of the form
  \[
    T \in \bigoplus_a D_a \otimes V_a
  \]
  where $D_a$ is a matrix and $V_a$ is a \emph{fusion-tree}, i.e. a clebsch-gordan tensor.
  Reviews and specific algorithms in the SU(2) case are given by \cite{schmollBenchmarkingGlobal$SU2$2020, schmollTensorNetworkSimulations2020, schmollProgrammingGuideTensor2020}.
  \cite[§~8]{schmollProgrammingGuideTensor2020} proposes that this method is appropriate for Anyonic simulation as well (although this is done earlier by singh \& pfeiffer.).
  Symmetries as categorified networks, as we will present in \cref{chap:?}, are explained in \cite{dissertation} 
  and \cite{devosTensorKitjlJuliaPackage2025} and in this work in \cref{sect:Anyons}.

  For a list of software implementing symmetries see \cref{app}.
\end{literature}


\begin{literature}[\bf{Anyons}]\label{Lit:Anyons}
  \emph{Anyons} as quasiparticles with exotic exchange statistics were originally propsed by \cite{wilczek1990} \cite{leinaasTheoryIdenticalParticles1977a}. Anyons are usually described as \emph{defects} in a (topologically orderd?) medium of dimension two (technically codimension after \cite{SS23AnyonsInTedK}).
  The exchange statistics of $N$ particles on this medium $M$ is goverend by the \emph{configuration space}
  \[
    \mathrm{Conf}_N(M) = \{(x_1,\ldots,x_N) \in M^N : x_i \neq x_j \} \setminus S_N.
  \]
  The statistics of the particle exchange is given by the representation of $\pi_1(\mathrm{Conf}_N(M))$ (the fundamental (homology?) group). If $M = \R^2$, this corresponds to the \emph{braid group} on $N$ strands $B_N$. If the dimension of $M \geq 3$ this corresponds to the symmetric group $S_N$. If these representations are unitary and the (ground state?) is degenerate, this allows for \emph{topological quantum computation} (c.f. \Cref{lit:tqc}).
\end{literature}



\begin{literature}[Anyonic Spin Chains and Tensor Networks]
  (Check if this is the first). A first appearance of \emph{Anyonic Tensor Networks} in terms of spin chains is given by \textcite{feiguinInteractingAnyonsTopological2007} as the \emph{golden chain}, which realizes a interacting $2$-local hamiltonian on a $n$ $\Fib$-anyon lattice. This is reviewed in \cite{trebstShortIntroductionFibonacci2008a} with generalizations to $3$-local and longer range hamiltonians. A general account outside of the $\Fib$-model can be found in \cite{gilsAnyonicQuantumSpin2013} and \cite{buicanAnyonicChainsTopological2017} (Check these again.). We can consider these approaches as a sort of manualization of the tensor network approach. Anyonic tensor networks are first introduced by \cite{pfeiferSimulationAnyonsTensor2010} and \cite{pfeiferSimulationAnyonsUsing2012}, developed out of the symmetry methods from \cref{sect:symmetricTensorNetworks} and both the \emph{isotopy} rules and anyonic operators defined in \cite{bondersonInterferometryNonAbelianAnyons2008}. Furthermore, stuctures such as an anyonic MERA are introduced. 

  We also need to discuss the topic of \emph{anyonic entanglement entropy}, as given in \cite{pfeiferMeasuresEntanglementNonAbelian2014,bondersonAnyonicEntanglementTopological2017,pfeiferSimulationAnyonsUsing2012}.
\end{literature}

\begin{literature}[Simulation of Anyons and String Nets]
  This should probably go in the previous section. \cite{liuMethodsSimulatingStringNet2022} \cite{kirchnerNumericalSimulationNonabelian2023}
\end{literature}

\begin{literature}[Anyons as Categories]\label{lit:anyonCat}
  The correspondence between the \emph{physical} treatment of anyons and the \emph{categorical treatment}
  of anyons as a \emph{braided fusion category} (modular functor?) was originally 
  descibed by \textcite[§~E, §~8]{kitaevAnyonsExactlySolved2006} based on previous work
  in quantum field theory already making use of modular categories and their diagrammatic
  calculus such as \cite{reshetikhinInvariants3manifoldsLink1991}. Further treatments are
  found in \cite{wang} and \cite{simon-nayak-das-sarma}. A modern treatment of the
  realization through the Configuration space of $n$-anyons if given by \cite{90eb98357acd4634abd18d906cc3811d}
  and in terms of modular functors by \cite{burtonShortGuideAnyons2016}.
\end{literature}

\begin{literature}[\bf Topological Quantum Computation]\label{lit:tqc}
  Fault tolerant \emph{Topological Quantum Computation (TQC)} using anabelian anyons was originally proposed by \textcite[§~7]{kitaevFaulttolerantQuantumComputation2003} (Originally published in 1997). Also \cite{freedmanNPQuantumField1998} and .... ??? has shown that there are nonabelian anyon models that are universal \cite{dawsonSolovayKitaevAlgorithm2005a} \cite{mochonAnyonsNonsolvableFinite2003} \cite{mochonAnyonComputersSmaller2004a}. 

  \cite{myersTopologicalQuantumGates2024a} gives nicer literature. \cite{simonTopologicalQuantum2023} is a great textbook account from one of the original authors on quantum compiling and computing. \cite{nayakNonAbelianAnyonsTopological2008b} is also one of the original accounts. This does not need to be any longer. Maybe mention categorical computation in general? but idk what htis is yet \cite{kongCategoricalComputation2023}.
\end{literature}

\begin{literature}[\bf Quantum Compiling]
  Some (Many?) anyon theories are universal for quantum computation. The \emph{Solovay-Kitaev theorem} \cite{kitaevQuantumComputationsAlgorithms1997}\cite{dawsonSolovayKitaevAlgorithm2005a} yields an efficient algorithm for approximating any quantum gate using only basic gates from a universal set. 
  Here we can consider 1.) The representation of braid circuits as a tensor nework operation and 2). A natural embedding of injection weaves as proposed by \textcite{simonTopologicalQuantumComputing2006a} where $2$ anyons are braided in their composite fusion channel by realizing this operation as a tensor reshape.
\end{literature}

\begin{literature}[Categorified Quantum Mechanics]
  There are multiple approach to a diagramatic and categorified approach to quantum mechanics. Generally one wants both spaces and maps represented by a framework that allows for parallel and entangled operations, i.e. has a tensor product structure. \cite{selingerSurveyGraphicalLanguages2011} gives a review of this \cite{vercleyenMathematicalStructureTensor2018} does for tensor networks specifically.

  TODO: find the XZ-calculus literature by Kissinger? Specific examples such as \emph{XZ}-calculus in the context of TQC is handled by \cite{ahmadiZXcalculusLanguageTopological2023}.

  TODO: How to mention that this results in or from $\Hilbfd$? Does it? The work from selinger et.al does not yet include anyons. For this, one needs \Cref{lit:anyonCat}.

  Also, take \cite[§~1]{satiEntanglementSectionsPushout2026} for a description
  of categorified mechanics. This stats that structure for \emph{entangled} processes
  is given by the monoidal category of complex vector spaces.
  
  % https://q.uiver.app/#q=WzAsNixbMSwwLCIoXFxNb2R7XFxDfSwgXFxvdGltZXMsIFxcQykiXSxbMSwyLCJcXE1vZHtcXEN9Il0sWzAsMCwiXFxzdGFja3JlbHtcXHRleHR7ZW50YW5nbGVkIHByb2Nlc3Nlc319e1xcdGV4dHtuby1jbG9uaW5nfX0iXSxbMCwyLCJcXHN0YWNrcmVse1xcdGV4dHtxdWFudHVtIHByb2Nlc3Nlc319e1xcdGV4dHtzdXBlcnBvc2l0aW9ucyBieX0gXFxvcGx1c30iXSxbMywwLCIoXFxIaWxiLFxcb3RpbWVzLD8pIl0sWzMsMiwiXFxIaWxiIl0sWzEsNSwiXFxkYWdnZXJcXHRleHR7LXN0cnVjdHVyZX0iXSxbMCw0LCJcXGRhZ2dlclxcdGV4dHstc3RydWN0dXJlfSJdLFsxLDBdLFs1LDRdXQ==
  \begin{tikzcd}
	  {\stackrel{\text{parallel quantum processes}}{\text{entanglement, no-cloning}}} & {(\Mod{\C}, \otimes, \C)} && {(\Hilb,\otimes,?)} \\
	  \\
	  {\stackrel{\text{quantum processes}}{\text{superpositions by} \oplus}} & {\Mod{\C}} && \Hilb
	  \arrow["{\dagger\text{-structure}}", from=1-2, to=1-4]
	  \arrow[from=3-2, to=1-2]
	  \arrow["{\dagger\text{-structure}}", from=3-2, to=3-4]
	  \arrow[from=3-4, to=1-4]
  \end{tikzcd}

  The categorical language specifically yield interpretations of theorems such as no-cloning and no deleting \cite{abramskyNoCloningCategoricalQuantum2012}.

  Additionally, many of the concepts of topological order are governed by category theory \cite{kongInvitationTopologicalOrders2022b}, as well as the topic of higher symmetries (should this be its own Lit sect if relevant? It is for MPO symmetries \cite{delcampHigherCategoricalSymmetries2024} \cite{kongAlgebraicHigherSymmetry2020}).
\end{literature}

\subsection{Tensors (Everything here is $fdVect_k$)}

Generally, we want to regard a tensor as a thing equipped with legs, where each legs holds a vector space. But we also want to be able to regard tensors as operators such as Hamiltonians.

Lets list some definitions of Tensor. A Tensor $T$ can be regarded as an element of a tensor product space
\[
  T \in \bigotimes_{j \in J} V_j
\]
where each $V_j$ is a finite-dimensional vector space and we say $T$ is a $\left\vert J \right\vert$-legged tensor.

Given !finite! dimensional Vector spaces $V$ and $W$, we have the identity that $V \otimes W^\ast \equiv Hom(W,V)$. This is called the \emph{hom-tensor duality}. (TODO: is this canonical? What about the other way around?)

In a basis expansion, this would relate the expansion of $V \otimes W^\ast$ as
\[
  T = \sum_{i}^{} v_i \otimes f_i
\]
where $f_i \in W^{\ast}   : W \to k$ with the map $T \in \Hom(W,V)$ as
\[
  T = (w \mapsto \sum_{i}^{} f_i(w) v_i).
\]


Now, if we have a Tensor $T \in V_1 \otimes \cdots \otimes V_n$, and since for any Vector space $V$, $(V^\ast)^\ast = V$, we can peel off any number of $V_{i}$s, take their duals and represent $T$ as a map
\[
  T : V_{n}^\ast \otimes \cdots \otimes V_{k+1} ^\ast \to V_1 \otimes \cdots \otimes V_k 
\]
Note that we present the domain here in a contravariant fashion due to the contravariance of the dual.

Lets look at an example, let $t \in V_1 \otimes V_2 \otimes W_{2}^\ast \otimes W_1^\ast$. We can represent as an operator in dirac notation:
\[
  t = \sum_{i_1,i_2,j_1,j_2}^{} T^{i_1 i_2}_{j_1 j_2}  \ket{v^1_{i_1}} \ket{v^2_{i_2}} \bra{w^1_{j_1}} \bra{w^2_{j_2}}
\] 
where $\ket{v_i}$ 

Graphically, we represent a general tensor
\[
  T : (W_{n}^{\ast} \otimes \cdots \otimes W_1) \to (V_1 \otimes V_2 \otimes \cdots \otimes V_n)
\]
as follows, where we choose incoming arrows for duals in the domain and outgoings for normal spaces, and on the bottom we choose incoming for duals and outgoing for normals.

